
\documentclass{article}
\usepackage{colortbl}
\usepackage{makecell}
\usepackage{multirow}
\usepackage{supertabular}

\begin{document}

\newcounter{utterance}

\centering \large Interaction Transcript for game `cladder', experiment `full\_v1.5\_default', episode 4304 with reflexion{-}Qwen3.5{-}4B.
\vspace{24pt}

{ \footnotesize  \setcounter{utterance}{1}
\setlength{\tabcolsep}{0pt}
\begin{supertabular}{c@{$\;$}|p{.15\linewidth}@{}p{.15\linewidth}p{.15\linewidth}p{.15\linewidth}p{.15\linewidth}p{.15\linewidth}}
    \# & \multicolumn{2}{c}{Player} && \multicolumn{2}{c}{Game Master} \\
    \hline

    \theutterance \stepcounter{utterance}  
    & & & \multicolumn{4}{p{0.6\linewidth}}{
        \cellcolor[rgb]{0.9,0.9,0.9}{
            \makecell[{{p{\linewidth}}}]{
                \texttt{\tiny{[P1$\langle$GM]}}
                \texttt{You are an expert in causal inference. The following question is not a typical commonsense query, but rather a meticulously designed question created by a professor specializing in causal inference, intended to assess the students' mastery of the course content.} \\
\texttt{Imagine a self{-}contained, hypothetical world with only the following conditions, and without any unmentioned factors or causal relationships: Pre{-}conditions has a direct effect on vaccination and lactose intolerance. Vaccination has a direct effect on lactose intolerance. The overall probability of getting the vaccine is 35\%. The probability of vaccine refusal and being lactose intolerant is 40\%. The probability of getting the vaccine and being lactose intolerant is 9\%. Is the chance of being lactose intolerant larger when observing getting the vaccine?} \\
\texttt{Start your answer with 'Yes' or 'No', followed by additional reasoning or evidence to support your explanation.} \\
            }
        }
    }
     \\ \\

    \theutterance \stepcounter{utterance}  
    & \multicolumn{4}{p{0.6\linewidth}}{
        \cellcolor[rgb]{0.561,0.737,0.561}{
            \makecell[{{p{\linewidth}}}]{
                \texttt{\tiny{[P1$\rangle$GM]}}
                \texttt{No} \\
\\ 
\texttt{**Reasoning:**} \\
\\ 
\texttt{To determine whether the chance of being lactose intolerant is larger when observing vaccination, we must calculate the conditional probability \$P(\textbackslash{}text\{Lactose Intolerance\} \textbackslash{}mid \textbackslash{}text\{Vaccinated\})\$ and compare it to the marginal probability \$P(\textbackslash{}text\{Lactose Intolerance\})\$.} \\
\\ 
\texttt{Let:} \\
\texttt{*   \$V\$ = Event of getting the vaccine.} \\
\texttt{*   \$L\$ = Event of being lactose intolerant.} \\
\\ 
\texttt{From the problem statement, we have the following probabilities:} \\
\texttt{1.  \$P(V) = 0.35\$} \\
\texttt{2.  \$P(\textbackslash{}neg V \textbackslash{}cap L) = 0.40\$ (Probability of vaccine refusal and being lactose intolerant)} \\
\texttt{3.  \$P(V \textbackslash{}cap L) = 0.09\$ (Probability of getting the vaccine and being lactose intolerant)} \\
\\ 
\texttt{**Step 1: Calculate the marginal probability of lactose intolerance, \$P(L)\$.**} \\
\texttt{The total probability of being lactose intolerant is the sum of the probabilities of being intolerant regardless of vaccination status:} \\
\texttt{\$\$P(L) = P(V \textbackslash{}cap L) + P(\textbackslash{}neg V \textbackslash{}cap L)\$\$} \\
\texttt{\$\$P(L) = 0.09 + 0.40 = 0.49\$\$} \\
\\ 
\texttt{**Step 2: Calculate the conditional probability of lactose intolerance given vaccination, \$P(L \textbackslash{}mid V)\$.**} \\
\texttt{Using the definition of conditional probability:} \\
\texttt{\$\$P(L \textbackslash{}mid V) = \textbackslash{}frac\{P(V \textbackslash{}cap L)\}\{P(V)\}\$\$} \\
\texttt{\$\$P(L \textbackslash{}mid V) = \textbackslash{}frac\{0.09\}\{0.35\} \textbackslash{}approx 0.257\$\$} \\
\\ 
\texttt{**Step 3: Compare the probabilities.**} \\
\texttt{*   Marginal probability \$P(L) = 0.49\$} \\
\texttt{*   Conditional probability \$P(L \textbackslash{}mid V) \textbackslash{}approx 0.257\$} \\
\\ 
\texttt{Since \$0.257 < 0.49\$, the probability of being lactose intolerant is actually **smaller** when observing that someone has received the vaccine. This indicates a negative association between vaccination and lactose intolerance in this specific hypothetical world, driven by the fact that the majority of lactose intolerant individuals (40\%) are vaccine refusers.} \\
\\ 
\texttt{Therefore, observing vaccination reduces the likelihood of lactose intolerance compared to the general population.} \\
            }
        }
    }
    & & \\ \\

    \theutterance \stepcounter{utterance}  
    & & & \multicolumn{2}{p{0.3\linewidth}}{
        \cellcolor[rgb]{0.9,0.9,0.9}{
            \makecell[{{p{\linewidth}}}]{
                \texttt{\tiny{[GM$|$GM]}}
                \texttt{No} \\
            }
        }
    }
    & & \\ \\

    \theutterance \stepcounter{utterance}  
    & & & \multicolumn{2}{p{0.3\linewidth}}{
        \cellcolor[rgb]{0.9,0.9,0.9}{
            \makecell[{{p{\linewidth}}}]{
                \texttt{\tiny{[GM$|$GM]}}
                \texttt{no} \\
            }
        }
    }
    & & \\ \\

    \theutterance \stepcounter{utterance}  
    & & & \multicolumn{2}{p{0.3\linewidth}}{
        \cellcolor[rgb]{0.9,0.9,0.9}{
            \makecell[{{p{\linewidth}}}]{
                \texttt{\tiny{[GM$|$GM]}}
                \texttt{game\_result = WIN} \\
            }
        }
    }
    & & \\ \\

\end{supertabular}
}

\end{document}
